\section{Conclusion}
\label{sec_concolusion}

%In this work, 
We propose {\papernameAbbrev}: a novel and interpretable framework for 3D visual grounding. By formulating the problem of 3D visual grounding from a natural language instruction as a sequence-to-sequence task, our approach predicts a chain of anchor objects that are subsequently utilized to localize the final target object. This sequential approach enhances interpretability and improves overall performance and data efficiency. Our framework is data-efficient and outperforms existing methods on the Nr3D and Sr3D datasets when trained on a limited amount of data. Furthermore, our proposed chain-of-thoughts module can easily be integrated into other architectures. Through extensive experiments, we demonstrate consistent performance gains over previous state-of-the-art methods operating on Referit3D.
Importantly, our approach does not rely on any additional manual annotations. Instead, we leverage automatic rule-based methods, syntactic parsing, and in-context learning to collect pseudo-labels for the anchor objects, thereby eliminating the laborious and time-consuming process of manually annotating anchors.
Overall, our work advances 3D visual grounding by making a step towards bridging the gap between machine perception and human-like understanding of 3D scenes.
